\subsection{So sánh tính năng giữa các nền tảng hiện có và hệ thống đề xuất}

\begin{table}[H]
\centering
\caption{So sánh tính năng giữa các nền tảng hiện có và hệ thống đề xuất}
\renewcommand{\arraystretch}{1.5}
\begin{tabular}{|p{5.5cm}|c|c|c|c|}
\hline
\textbf{Nhóm tính năng} & \textbf{CellphoneS} & \textbf{FPT Shop} & \textbf{Thế Giới Di Động} & \textbf{Hệ thống đề xuất} \\
\hline
Tìm kiếm và lọc sản phẩm & Có & Có & Có & Có \\
\hline
Phân loại và hiển thị sản phẩm & Có & Có & Có & Có \\
\hline
Đánh giá và phản hồi sản phẩm & Có & Có & Có & Có \\
\hline
Quản lý giỏ hàng và đơn hàng & Có & Có & Có & Có \\
\hline
Hệ thống khuyến mãi và ưu đãi & Có & Có & Có & Có \\
\hline
Giao diện thân thiện, tối ưu UX/UI & Có & Có & Có & Có \\
\hline
Tư vấn sản phẩm & Không & Không & Hạn chế & Thông minh \\
\hline
\end{tabular}
\renewcommand{\arraystretch}{1.0}
\end{table}

Từ kết quả khảo sát ba nền tảng thương mại điện tử lớn tại Việt Nam, có thể nhận thấy rằng mỗi hệ thống đều có định hướng riêng trong việc phát triển và phục vụ người dùng.

Tuy nhiên, nhóm nhận thấy rằng mặc dù các doanh nghiệp trong ngành đã có bước tiến lớn trong việc số hóa hoạt động bán lẻ, các nền tảng thương mại điện tử trong mảng bán lẻ điện thoại di động và các đồ dùng công nghệ tại Việt Nam chưa thật sự tích hợp Trợ lý Ảo thông minh có khả năng giao tiếp tự nhiên và duy trì luồng hỗ trợ nhiều lượt và phi tuyến tính.

Khoảng trống này mở ra cơ hội để nhóm nghiên cứu và phát triển một mô hình thương mại điện tử tích hợp Trợ lý Ảo AI Agent, hướng đến việc nâng cao trải nghiệm người dùng, tự động hóa quy trình giao dịch, và tăng tính cá nhân hóa trong tương tác mua sắm.